\documentclass{article}

\usepackage[margin=1in]{geometry}
\usepackage{amsmath,amsthm,amssymb}
\usepackage[spanish]{babel}
\usepackage{enumitem}
\usepackage{pgfplots}
\pgfplotsset{compat=1.15} % Asegura compatibilidad de pgfplots
\usepgfplotslibrary{fillbetween}

\theoremstyle{definition}
\newtheorem{definition}{Definición}[section]

% Number sets
\newcommand{\R}{\mathbb{R}}
\newcommand{\Z}{\mathbb{Z}}
\newcommand{\N}{\mathbb{N}}
\newcommand{\Q}{\mathbb{Q}}
\newcommand{\C}{\mathbb{C}}

% Topology operators
\newcommand{\cl}[1]{\overline{#1}}              % Closure
\newcommand{\Int}[1]{{#1}^\circ}                % Interior
\newcommand{\bd}{\partial}                       % Boundary
\newcommand{\Ent}[1]{\mathcal{O}(#1)}           % Neighborhoods (Entornos)
\newcommand{\pf}{\mathcal{P}_F}                 % Finite parts
\newcommand{\partes}{\mathcal{P}}               % Power set

% Useful shortcuts
\newcommand{\sii}{\iff}                          % Si y sólo si
\newcommand{\imp}{\implies}                   % Implies
\newcommand{\setdiff}{\smallsetminus}            % Set difference

\newcommand{\coint}[2]{\left[ #1, #2 \right)}

% Exercise environment
\newenvironment{ejercicio}[1]{\begin{trivlist}
\item[\hskip \labelsep {\bfseries Ejercicio}\hskip \labelsep {\bfseries #1.}]}{\end{trivlist}}

\setlist[enumerate,1]{label=(\alph*), leftmargin=*, itemsep=2pt}

\begin{document}

\title{Topología: Práctica III}
\author{Bustos Jordi\\Espacios Conexos}

\maketitle

\begin{ejercicio}{1}
    Sean \((X, \tau)\) un ET, \( \{ A_\alpha \} \) una colección de subespacios conexos de \(X\) y \(A\) un subespacio conexo de \(X\). Demostrar que si \( A \cap A_\alpha \neq \varnothing \) para todo \( \alpha \), entonces \(A \cup \left( \bigcup_\alpha A_\alpha \right)\) es conexo.
\end{ejercicio}
\begin{proof}
    Es claro que \( A \, \cup \, \bigcup_\alpha A_\alpha = \bigcup_\alpha (A \cup A_\alpha) \).
    Dado que \( A \) y \( A_\alpha \) son conexos y \( A \cap A_\alpha \neq \varnothing \), entonces \( A \cup A_\alpha \) es conexo para cada \( \alpha \). Aplicando el mismo razonamiento \( \bigcup_\alpha (A \cup A_\alpha) \) es conexo.
\end{proof}

\begin{ejercicio}{2}
    Probar que si \( X \) es un conjunto infinito, entonces \(X\) es conexo con la topología de los complementos finitos (o topología finita).
\end{ejercicio}
\begin{proof}
    En efecto, supongamos que tenemos una \( X \)-separación de \( X \) tal que \( U, V \) abiertos de la topología finita, \( U \cup V = X \) y
    \( U \cap V \cap X = U \cap V = \varnothing \) y veamos que no puede ocurrir que \( U \cap X \neq \varnothing \) y \( V \cap X \neq \varnothing \).
    En efecto, como los complementos de los abiertos son finitos, entonces \( |U^c| = n \in \N \) y \( |V^c| = m \in \N \).
    Luego, existen menos de \( m \) elementos en \( U \) que no están en \( V \) pues \( U \cup V = X \).
    De la misma manera existen menos de \( n \) elementos en \( V \) que no están en \( U \) y \( U \cap V \) puede dejar afuera a lo sumo \( n + m \) elementos de \( X \), pero como \( X \) es infinito y \( U \cap V = \varnothing \) debe ser \( U = \varnothing \) o \( V = \varnothing \).
\end{proof}

\begin{ejercicio}{3}
    Sean \(\tau_1\) y \(\tau_2\) dos topologías en \(X\) tal que \(\tau_1 \subseteq \tau_2\). ¿Qué puede decirse de la conectividad de \(X\) en una de sus topologías respecto de la conectividad en la otra?
\end{ejercicio}
\begin{proof}
    Sea \( \tau_1 = \{ \varnothing, X \} \) y \( \tau_2 = \{ \varnothing, A, A^c, X \} \) con \( A \subset X \) es claro que \( \tau_1 \subseteq \tau_2 \).
    Además \( \tau_1 \) es conexa pero \( \tau_2 \) no pues \( A \) es clopen no trivial de \( \tau_2 \).

    Por otro lado, si vale que si \( X \) es conexo con \( \tau_2 \) entonces \( X \) es conexo con \( \tau_1 \) si \( \tau_1 \subseteq \tau_2 \). En efecto,
    dados \( U, V \in \tau_1 \) con \( U \cup V = X \) y \( U \cap V \cap X = \varnothing \), la separación no puede ser fuerte pues \( U, V \in \tau_2 \) y eso implicaría que \( X \) no es conexo en \( \tau_2 \).
\end{proof}

\begin{ejercicio}{4}
    Sea \(A\) un subespacio conexo de \(X\). Analizar si \( \Int{A} \) y \(\partial A\) son conexos. ¿Vale la recíproca?
\end{ejercicio}
\begin{proof}
    Para ver \( \partial A \) no es conexo basta con considerar \( A \) un anillo en \( \R^2 \) y su frontera \( \partial A \) es la unión de dos círculos concéntricos, que no son conexos entre sí.
    Para ver que \( \Int{A} \) no es conexo basta con considerar dos discos tangentes de radio \( 1 \) y \( A \) la unión de ambos. Entonces \( A \)
    es conexo pero \( \Int{A} \) es la unión de dos discos abiertos tangentes, que no son conexos entre sí pues se puede tomar cada disco para hacer la separación.

    Veamos si \( \Int{A} \) y \( \partial A \) es conexo, entonces \( A = \Int{A} \cup \partial A \) es conexo. Recordemos que la definición de frontera es \[
        \partial A = \cl{A} - \Int{A} = \cl{A} - \Int{A} = \cl{A} - \Int{A}
    \]
    Consideremos el siguiente contraejemplo: \( A = \{ (x, y) \in \R^2 : x^2 + y^2 < 1 \} \cup \{ (x, 0) \in \Q^2 : 1 \leq x \leq 2 \} \).
    Entonces \[ \cl{A} = \{ (x, y) \in \R^2 : x^2 + y^2 = 1 \} \cup \{ (x, 0) \in \R^2 : 1 \leq x \leq 2 \} \] y \[ \Int{A} = \{ (x, y) \in \R^2 : x^2 + y^2 < 1 \} \] que ambos son conexos. Sin embargo,
    \( A \) no es conexo pues podemos considerar el abierto \( U = \{ (x, y) \in \R^2 : x > \sqrt{2} \} \) y \( V = \{ (x, y) \in \R^2 : x < \sqrt{2} \} \) que forman una separación de \( A \).
\end{proof}

\begin{ejercicio}{5}
    Dar un contraejemplo para cada una de las siguientes afirmaciones falsas:
    \begin{enumerate}
        \item Un conjunto conexo implica que es localmente conexo.
        \item Un conjunto arcoconexo es localmente arcoconexo.
        \item La intersección de dos conjuntos conexos es un conjunto conexo.
        \item La intersección de 2 conjuntos localmente conexos es un conjunto localmente conexo.
    \end{enumerate}
\end{ejercicio}
\begin{proof}
    Para el primer punto consideremos el conjunto en \( \R^2 \) dado por \[ A = \{ (t, \sin(1/t)) : 0 < t \leq 1 \} \cup \{ (0, y) : -1 \leq y \leq 1 \} \] Por la \textbf{Observación 3.2.6 (iii)} sabemos que \( A \) es conexo
    pero no es localmente conexo pues cualquier entorno de \( (0, y) \) intersecta a \( A \) en un conjunto que no es conexo.

    Para el segundo punto consideremos el mismo conjunto con la siguiente modificación \[ A = \{ (t, \sin(1/t)) : 0 < t \leq 1 \} \cup \{ (0, y) : -1 \leq y \leq 1 \} \cup \{ (t, -1) : 0 \leq t \leq 1 \} \]
    Entonces \( A \) es arcoconexo, pero con el mismo razonamiento que el punto anterior \( A \) no es localmente arcoconexo.

    Para el tercer punto, sin salir de \( \R^2 \), podemos considerar \( A = \{ (x, y) \in \R^2 : 4 \leq y \leq 5 \} \) y
    \( B = \{ (x, y) \in \R^2 : x^2 \leq y \leq x^2 + 1 \} \), es fácil ver que la intersección son dos conjuntos disjuntos.

    Finalmente, para el cuarto punto, consideremos la siguiente modificación del seno topológico:\begin{equation*}
        f(x) = \begin{cases}
            \sin(1/x) & \text{si } x > 0    \\
            0         & \text{si } x \leq 0
        \end{cases}
    \end{equation*}
    y consideremos \( A = \{ (x, y) \in \R^2 : y \leq f(x) \} \) y \( B \) el semiplano superior. Entonces \( A \) y \( B \) son localmente conexos y su intersección es \[
        A \cap B = \{ (x, y) \in \R^2 : 0 \leq y \leq f(x) \} \cup \{ (x, y) \in \R^2 : x \leq y \leq 0 \}
    \]

    Si analizamos un entorno de \( (0, 0) \) tomando \( \varepsilon > 0 \) el entorno relativo \( A \cap B \cap B_{\varepsilon}((0, 0)) \) contiene al origen y a una cantidad infinita de islas del seno topológico relleno, como las islas están separadas
    por espacios vacíos, entonces \( A \cap B \cap B_{\varepsilon}((0, 0)) \) no es conexo y \( A \cap B \) no es localmente conexo.
\end{proof}

\begin{tikzpicture}
    \begin{axis}[
            axis lines=middle,
            xlabel={\( x \)},
            ylabel={\( y \)},
            xmin=-0.4, xmax=1.2,
            ymin=-1.0, ymax=1.0,
            samples=600,
            domain=0.015:1.2,
            enlargelimits=false,
            axis on top,
            width=15cm,
            height=9cm,
            xtick=\empty,
            ytick=\empty,
            grid=major,
            grid style={gray!20, thin}
        ]

        % 1. Semiplano superior B (fondo suave)
        \fill[blue!8] (axis cs:-0.4, 0) rectangle (axis cs:1.2, 1.0);
        \draw[blue!40, thin] (axis cs:-0.4, 0) -- (axis cs:1.2, 0);
        \node[blue!70, anchor=south west, font=\small] at (axis cs:0.5, 0.5) {\( B = \{y \ge 0\} \)};

        % 2. Conjunto A
        \addplot [name path=F, draw=none] {x * (sin(deg(15/x)) - 0.5)};
        \addplot [name path=Fneg, draw=none, domain=-0.4:0] {x};
        \addplot [name path=bot, draw=none, domain=-0.4:1.2] {-1};

        \addplot [red!12, forget plot] fill between [of=F and bot];
        \addplot [red!12, forget plot] fill between [of=Fneg and bot];
        \fill[red!12] (axis cs:0,0) -- (axis cs:0.015, {0.015 * (sin(deg(15/0.015)) - 0.5)}) -- (axis cs:0.015, -1) -- (axis cs:0,-1) -- cycle;

        % 3. Intersección A ∩ B
        \begin{scope}
            \clip (axis cs:-0.4, 0) rectangle (axis cs:1.2, 1.0);

            \addplot [teal!60!black, forget plot] fill between [of=F and bot];
            \addplot [teal!60!black, forget plot] fill between [of=Fneg and bot];
            \fill[teal!60!black] (axis cs:0,0) -- (axis cs:0.015, {0.015 * (sin(deg(15/0.015)) - 0.5)}) -- (axis cs:0.015, -1) -- (axis cs:0,-1) -- cycle;
        \end{scope}

        % 4. Frontera de A (curva f(x))
        \addplot [red!80!black, very thick, forget plot] {x * (sin(deg(15/x)) - 0.5)};
        \addplot [red!80!black, very thick, domain=-0.4:0, forget plot] {x};
        \draw[red!80!black, very thick] (axis cs:0,0) -- (axis cs:0.015, {0.015 * (sin(deg(15/0.015)) - 0.5)});

        % 5. Círculo de entorno alrededor del origen (punteado)
        \draw[dashed, teal!40, line width=1pt] (axis cs:0,0) circle (0.12cm and 0.12cm);

        % 6. Puntos clave y etiquetas
        \fill[teal!80] (axis cs:0,0) circle (2.5pt);
        \node[teal!80, anchor=south east, font=\small, inner sep=2pt] at (axis cs:-0.02, -0.05) {\( (0,0) \)};

        % Etiqueta del conjunto A
        \node[red!70, anchor=north west, font=\small] at (axis cs:0.1, -0.7) {\(A = \{y \le f(x)\}\)};

        % Etiqueta descriptiva de las islas
        \node[teal!70, anchor=south west, font=\small, align=left] at (axis cs:0.22, 0.35) {\textbf{Islas} \\ \textbf{de} \( A \cap B \)};
    \end{axis}
\end{tikzpicture}

\clearpage

\begin{ejercicio}{6}
    Sean \((X, \tau)\) un ET y \(A, B \subset X\) tales que \(A \cup B\) y \(A \cap B\) son conexos.
    \begin{enumerate}
        \item Si \( A \) y \( B \) son abiertos entonces  \( A \) y \( B \) son conexos.
        \item Si \( A \) y \( B \) son cerrados entonces \( A \) y \( B \) son conexos.
    \end{enumerate}
\end{ejercicio}
\begin{proof}
    Supongamos que \( A \) y \( B \) son abiertos tal que \( A \cap B \neq \varnothing \) y, s.p.d.g, que \( A \) no es conexo.
    Entonces existe una \( X \)-separación fuerte \( E \), \( F \) de \( A \) tal que
    \( A \subseteq E \cup F \), \( A \cap E \cap F = \varnothing \), \( A \cap E \neq \varnothing \) y
    \( A \cap F \neq \varnothing \).

    Tenemos que \( A \subseteq E \cup F \implies A \cap B \subseteq (E \cap B) \cup (F \cap B) \) y \( E \cap B, F \cap B \) son abiertos pues \( E, F, B \) lo son.

    \( A \cap E \cap F = \varnothing \implies A \cap E \cap F \cap B = \varnothing \equiv (A \cap B) \cap (E \cap B) \cap (F \cap B) = \varnothing \). Como \( A \cap B \) es conexo, vale que
    \( A \cap B \cap E = \varnothing \) o \( A \cap B \cap F = \varnothing \). Supongamos que \( A \cap B \cap E = \varnothing \), como \( A \cap E \neq \varnothing \)
    y \( A \cap B \neq \varnothing \implies E \cap B \neq \varnothing \). Luego, \( A \cup B \subseteq (E \cup F) \cup B \equiv A \cup B \subseteq E \cup (F \cup B) \).

    Se tiene que \( E \cap F \cap B = \varnothing \), \( ( A \cup B ) \cap F = (A \cap F) \cup (B \cap F) \neq \varnothing \) y \( (A \cup B) \cap E \neq \varnothing \). Como \( E \), \( F \cup B \) son abiertos,
    \( (A \cup B) \cap E \cap (F \cup B) = \varnothing \) y \( (A \cup B) \cap E \neq \varnothing \) y \( (A \cup B) \cap F \neq \varnothing \). Absurdo! Pues \( A \cup B \) es conexo.

    La otra implicación sale considerando que \( A^c \) y \( B^c \) son abiertos y aplicando el mismo razonamiento al caso anterior.
\end{proof}

\begin{ejercicio}{7}
    En \( \N \) se define \( \tau = \{ I_n : n \in \N \} \cup \{ \varnothing, \N \} \), con \( I_n = \{ 1, 2, \ldots, n \} \). ¿Cuáles son los subconjuntos conexos?
\end{ejercicio}
\begin{proof}
    Es claro que todo subconjunto \( A \) es conexo pues si hubiese una separación con dos abiertos de la topología \( I_n \), \( I_m \), como vale que \[
        I_n \subseteq I_m \text{ o } I_m \subseteq I_n
    \]
    no se puede hacer una separación de \( A \) con \( I_n \) e \( I_m \) pues \( A \subseteq I_n \cup I_m \equiv A \subseteq I_n \text{ o } A \subseteq I_m \) y, si \( n \leq m \), \( A \cap I_n \cap I_m = A \cap I_n = A \neq \varnothing \).
\end{proof}

\begin{ejercicio}{8}
    Sea \((X, \tau)\) un espacio topológico. Sea \( x \in U \subseteq X \) con \(U\) una componente conexa de \(X\). Demostrar que
    \begin{enumerate}
        \item \(U\) es el mayor subconjunto conexo de \(X\) que contiene a \(x\).
        \item Toda componente conexa es cerrada.
    \end{enumerate}
\end{ejercicio}
\begin{proof}
    Consideremos \( U \) la componente conexa de \( X \) que contiene a \( x \) y supongamos que \( V \) es un conexo tal que \( x \in V \) y
    \( U \subset V \). Entonces, como \( V \) es estrictamente más grande que \( U \), existe \( y \in V \setminus U \), pero como existe
    un conexo (\(V\)) tal que \( x, y \in V \), entonces \( x \overset{con}{\sim} y \) y, por lo tanto, \( y \in U \). Absurdo! Por lo tanto, \( U \) es el mayor subconjunto conexo de \( X \) que contiene a \( x \).

    Veamos ahora que toda componente conexa \( U \) es cerrada. Por el absurdo, supongamos que no lo es. Entonces, existe \( y \in \cl{U} \setminus U \). Consideremos ahora
    el conjunto \( C = U \cup \{ y \} \). Si \( C \) admite una \( X \)-separación fuerte \( E, F \in \tau \) tal que \( C \subseteq E \cup F \),
    \( C \cap E \cap F = \varnothing \) y \( C \cap E \neq \varnothing \) y \( C \cap F \neq \varnothing \), entonces como \( U \subseteq C \subseteq E \cup F \),
    sin pérdida de generalidad podemos suponer que \( U \subseteq E \) y \( C \cap F = (U \cup \{ y \}) \cap F = \{ y \} \).
    Pero notemos ahora que \( \forall V \in \mathcal{O}^a(y) \), \( V \cap (U \setminus \{ y \}) = V \cap U \neq \varnothing \), tomando \( F = V \) en la definición llegamos a que \[
        F \cap (U \setminus \{ y \}) = F \cap U = \varnothing
    \]
    lo cual contradice que \( y \in \cl{U} \). Por lo tanto, \( C \) es conexo y \( U \subset C \) que contradice la maximalidad de \( U \). Por lo tanto, \( y \in U \) y \( U \) es cerrado
\end{proof}

\begin{ejercicio}{9}
    Mostrar que en un espacio localmente arcoconexo las componentes conexas y las componentes arcoconexas son iguales.
\end{ejercicio}
\begin{proof}
    Sea \( x \in X \). Si llamamos \( E \) y \( F \) a la componente conexa y arcoconexa respectivamente, entonces \( F \subseteq E \) es trivial pues todo conjunto arco-conexo es conexo y \( E \) es el más grande.
    Veamos ahora que \( E \subseteq F \), para eso, recordemos que como \( E \) es conexo vale que los únicos clopen relativos a \( E \) son \( \varnothing \) y \( E \).
    Sea \( y \in F \), \( \exists U \in \beta_x \subseteq \mathcal{O}^a(z) \) tal que \( y \in U \subseteq F \). Como \( U \cup F \) es arco-conexo, pues tienen un punto en común, entonces \( U \subseteq F \) por ser
    \( F \) maximal. Luego, para todo \( y \in F \) puedo encontrar un entorno abierto \( U \) contenido en \( F \) y por lo tanto \( F \) es abierto. Sea \( F^c \) la unión de todas las componentes arco-conexas distintas de \( F \),
    como todas son abiertas, \( F^c \) es abierto y \( F \) es cerrado. Luego, \( F \) es clopen en \( X \) y si consideramos el subespacio \( E \) con la topología relativa, \( F \) es clopen en \( E \).
    Por lo tanto, como \( F \cap E \neq \varnothing \) se sigue que \( F \cap E = E \implies E \subseteq F \).
\end{proof}

\begin{ejercicio}{10}
    Sea \(X\) un espacio localmente arcoconexo. Demuestre que cada conjunto abierto y conexo en \(X\) es también arcoconexo.
\end{ejercicio}
\begin{proof}
    Sea \( U \) un conjunto abierto y conexo en \( X \). Para cada \( x \in U \) consideremos \( C_x = \{ y \in U : \overset{con}{\sim} x \} \) la componente arcoconexa de \( x \) en \( U \) y, si logramos demostrar que es clopen relativo a \( U \), entonces \( C_x = U \) y, por lo tanto, \( U \) sería arcoconexo.

    Veamos primero que es abierto: En efecto, dado \( y \in C_x \) veamos que existe un entorno \( V \) de \( y \)
    completamente contenido en \( C_x \). Para eso, como \( X \) es localmente arcoconexo, existe un entorno arcoconexo \( V \subseteq U \) de \( y \).
    Además, dado \( z \in V \) es claro que puedo conectar \( y \) con \( z \) mediante un arco contenido en \( V \) y conectar
    \( x \) con \( y \) mediante un arco contenido en \( C_x \), entonces \( x \) y \( z \) están conectados por un arco contenido en \( C_x \cup V \), pero entonces \( z \in C_x \) y, luego \( V \subseteq C_x \).
    Esto nos dice que \( C_x \) es abierto.

    Veamos ahora que es cerrado: En efecto, dado \( y \in U \setminus C_x \) y \( V \subseteq U \) el entorno arcoconexo de la base de entornos de \( y \).
    Supongamos que \( V \cap C_x \neq \varnothing \), entonces existe \( z \in V \cap C_x \) y, como \( V \) es arcoconexo,
    puedo conectar a \( z \) con \( y \) y a \( y \) con \( x \) mediante un arco contenido en \( V \cup C_x \), esto implica que \( y \in C_x \), absurdo! Por lo tanto,
    \( V \cap C_x = \varnothing \) y \( U \setminus C_x \) es abierto. Por lo tanto, \( C_x \) es cerrado.

    Luego, \( C_x \) es clopen relativo a \( U \) conexo \( \therefore U = C_x \) y \( U \) es arcoconexo.
\end{proof}

\begin{ejercicio}{11}
    Mostrar que un conjunto infinito dotado de la topología co-finita es conexo.
\end{ejercicio}
\begin{proof}
    Sea \( U \in \tau \) supongamos que tenemos una separación de \( U \). Luego, \( \exists \, E, F \in \tau \) tal que
    \( U \subseteq E \cup F \), \( U \cap E \cap F = \varnothing \). Sea \( U^c = {\{ x_i \}}_{i=1}^n \) se tiene que \( E \cap F \subseteq U^c \) pero \( E \cap F \in \tau \) por lo que
    debe ser \( E \cap F = \varnothing \) y, por lo tanto aplicando el mismo razonamiento que en el \textbf{Ejercicio 2}, \( E = \varnothing \) o \( F = \varnothing \). Por lo que no
    puede haber una separación de \( U \) y \( U \) es conexo.
\end{proof}

\begin{ejercicio}{12}
    Sean \(x\) e \(y\) dos elementos distintos en la misma componente conexa de \(X\). Demostrar que cualquier clopen de \(X\) tiene ambos o a ninguno.
\end{ejercicio}
\begin{proof}
    Dados \( x \), \( y \) en la misma componente conexa \( C \). Supongamos que \( U \) es un clopen de \( X \) tal que \( x \in U \) y \( y \notin U \). Entonces \( V = X \setminus U \) es un clopen de \( X \) tal que 
    \( y \in V \) y \( x \notin V \). Luego, \( U \cup V = X \) y \( U \cap V = \varnothing \) es una separación de \( X \) tal que \( x \in U \) e \( y \in V \).
    Ahora notemos que \( C_1 = U \cap C \) y \( C_2 = V \cap C \) son abiertos relativos a \( C \) tales que \( C_1 \cup C_2 = C \), 
    \( \{ x \} \subseteq C_1 \cap C \) y \( \{ y \} \subseteq C_2 \cap C \) y \( C_1 \cap C_2 = \varnothing \). Esto contradice que \( C \) es conexo.
    Por lo tanto, no puede existir un clopen de \( X \) que contenga a \( x \), pero no a \( y \) o que contenga a \( y \) pero no a \( x \). Por lo tanto, cualquier clopen de \( X \) tiene ambos o a ninguno. 
\end{proof}

\clearpage

\subsection*{Ejercicios Complementarios}

\begin{ejercicio}{C1}
    Sea \(Y \subseteq X\) ambos conexos. Si \((U, V)\) es una separación fuerte de \(Y^c\) entonces \(Y \cup U\) y \(Y \cup V\) son conexos.
\end{ejercicio}

\begin{proof}
\end{proof}

\begin{ejercicio}{C2}
    Si \(Y\) es un subespacio conexo de \(X\) tal que \(Y \cap A \neq \varnothing\) y \(Y \cap A^c \neq \varnothing\) entonces \(Y \cap \partial A \neq \varnothing\).
\end{ejercicio}

\begin{proof}
\end{proof}

\begin{ejercicio}{C3}
    Considerar \(F_1 \supseteq F_2 \supseteq \cdots \supseteq F_n \supseteq \cdots\) una sucesión infinita decreciente de conjuntos cerrados y conexos en \(\R^2\) (con la topología usual). ¿Es \(\bigcap_{n \geq 1} F_n\) conexo?
\end{ejercicio}

\begin{proof}
\end{proof}

\begin{ejercicio}{C4}
    Probar que la imagen de un conjunto arcoconexo bajo una función continua es un conjunto arcoconexo.
\end{ejercicio}

\begin{proof}
\end{proof}

\begin{ejercicio}{C5}
    Considerar \(X = [0, 1] \times [0, 1]\) dotado con el orden lexicográfico y la topología de orden respectiva. Mostrar que \(X\) es conexo pero no arcoconexo.
\end{ejercicio}

\begin{proof}
\end{proof}

\end{document}